% Compile from this directory with: pdflatex project_16.tex (twice).
\documentclass[11pt,a4paper]{article}
\usepackage[T1]{fontenc}
\usepackage[utf8]{inputenc}
\usepackage[margin=22mm]{geometry}
\usepackage{lmodern}
\usepackage{amsmath,amssymb}
\usepackage{enumitem}
\usepackage{xcolor}
\usepackage{microtype}
\usepackage{hyperref}
\definecolor{epflred}{HTML}{E3000B}
\hypersetup{colorlinks=true,urlcolor=epflred,linkcolor=epflred,
  pdftitle={ENG-654 Project 16: Orientation-Dependent Connectivity on custom\_6R}}
\setlength{\parindent}{0pt}
\setlength{\parskip}{6pt}
\setlist[itemize]{leftmargin=*,itemsep=5pt,topsep=4pt}
\urlstyle{tt}

\begin{document}
{\small\textbf{EPFL \textbar\ ENG-654}\hfill Kinematics-Grounded Motion Planning for Robots}
\vspace{5mm}

{\color{epflred}\Large\bfseries Project 16}

{\LARGE\bfseries Orientation-Dependent Connectivity on custom\_6R\par}

\textbf{Format:} Group of two students; simulation-based project.\\
\textbf{Course foundations:} Lectures 4--7; Exercises 2--3.

\section*{Motivation}
The visible spatial trace of a tool does not fully describe its task.
For the same wrist-center curve, different orientation schedules can create
different wrist singularities and feasible joint-space continuations. This
project studies that dependence while holding the positioning problem fixed.

\section*{Task}
\textbf{Overall objectives.} Compare a small set of orientation schedules on one custom\_6R
wrist-center curve and explain how they alter the feasible continuations of
its arm and wrist IK solutions.

\begin{itemize}
  \item \textbf{Validate the robot and the chosen curve.} Use the exact
custom\_6R model, with the custom\_3R positioning arm and spherical wrist.
Check the D--H/URDF correspondence and fixed tool transform. Choose one
short closed wrist-center curve from a supplied feasible-path example.

  \item \textbf{Map the common arm problem.} Plot the positioning determinant and
zero contours in $(q_2,q_3)$ and map the two-/four-arm-IK workspace regions.
Enumerate the initial arm IKs and continue the arm candidates over the
curve. Overlay these trajectories so that all orientation experiments share
one clearly documented positioning baseline.

  \item \textbf{Define three orientation schedules.} Compare a constant tool
orientation, that orientation composed with one fixed offset rotation, and
a smooth periodic orientation variation. For example use
$R(s)=R_0\exp(\beta\sin(2\pi s)\widehat{u})$ with fixed unit axis $u$.
Record the angle and derive tool position consistently from the wrist center.

  \item \textbf{Lift and classify the full-pose candidates.} Complete each
continued arm candidate with the wrist solutions and track each starting
flip separately. Compare full-pose residuals, wrist/arm regularity, joint
limits, and endpoint configuration. Do not interpret an orientation change
as permission to jump between unrelated IKs.

  \item \textbf{Explain changes in connectivity.} Build a compact chart of path
progress versus candidate identity for each schedule, marking feasible
segments and failure reasons. Compare where a complete arm continuation
loses a full-pose lift. Refine one boundary and explain it using the
corresponding wrist singularity factor or active joint limit.
\end{itemize}

\newpage
\section*{Specifics and scope}
Use one wrist-center curve, three orientation schedules, and all
recovered regular initial candidates (up to eight). Begin with approximately
150 samples. The course arm/wrist IK and a reference curve are available for
reuse. Investigate the prescribed schedules; optimizing over arbitrary
orientation functions is outside scope.

Check both translation and orientation residuals
for every claimed full-pose IK. Document angle periodicity and limits,
Jacobian conventions, and any scaling used for singular values. State all
fixed coordinates in reduced plots; a two-dimensional slice does not show
the entire 6R singular set. Refine decisive transitions, and distinguish
sampled evidence from a continuous-path guarantee. Collision freedom is
not implied by these kinematic checks.

\section*{Expected outcome}
Explain why the same positioning task can have different full-pose
feasibility depending on orientation. Deliver the common arm maps, three
explicit pose schedules, candidate feasibility charts, and a comparison of
the global joint-space continuations.

Submit reproducible code, model parameters and test data, the required plots,
a concise 5--6-page report, and a short simulation demonstration. Explain
both successful cases and failures; conclusions must follow the numerical
and geometric evidence rather than an expected outcome.

\section*{Resources}
The custom\_6R consists of the custom\_3R positioning arm plus a
spherical wrist. The arm is orthogonal with non-intersecting consecutive axes.
Its URDF is \path{lectures_main/assets/models/custom_6R/custom_6R_new.urdf};
the corresponding STL meshes are in that directory. Inspect wrist-axis
concurrency and base/tool transforms using the lecture website.
Use Lectures 4--6 for the custom arm and wrist structure, and Lecture 7 for
continuation. The supplied reference curve must be interpreted in the exact
custom\_6R wrist-center frame. Existing position-only visualizations must be
complemented with full-pose validation.

Use the lecture website to verify D--H parameters, visualize IKs and
singularities, and simulate paths. All listed paths are relative to the
repository root. Start the website following \path{lectures_main/README.md}.

\section*{Workload and collaboration}
Budget approximately 60 hours per student (120 person-hours per pair)
for this 2-ECTS project, following EPFL's 30-hours-per-credit convention.\footnote{\href{https://www.epfl.ch/education/teaching/teaching-guide-2/getting-started/design-a-course_1/}{EPFL: Design a Course, workload guidance}.}
Deduct any lectures or exercises included in those credits. Allocate roughly
20 team-hours to model validation, 50 to the core work, 30 to experiments,
and 20 to reporting. Reuse the specified IK and simulations. Share validation
and interpretation even if modelling and implementation are divided between
students. Hardware execution is outside scope.

\end{document}
